\vspace{-1mm}
\section{Introduction}
\vspace{-2mm}
% 方法主图
\begin{figure*}[t]
    \centering
    \includegraphics[width=0.85\textwidth]{figs/intro_main.pdf}
    \vspace{-2mm}
    \caption{
    From External Control (Top) to State-of-Thought (Bottom).
    }
    \vspace{-7mm}
    \label{fig:intro_main}
\end{figure*}

Reasoning in Large Language Models (LLMs) is fundamentally the process of organizing historical evidence context for subsequent decisions
\citep{press2023measuring,trivedi2023interleaving,yao2023react}.
From this perspective, progress in test-time compute comes from better strategies for organizing the evidence context with additional computation, thereby improving reasoning capabilities \citep{brown2024large,snell2024scaling}.
One dominant line relies on externally structured reasoning programs, requiring the model to follow pre-specified reasoning formats such as step-by-step \citep{wei2022chain}, planning \citep{wang2023plan}, or revision \citep{madaan2023self}.
Another line allocates more test-time search, exploring additional samples \citep{brown2024large,wang2023selfconsistency}, branches \citep{yao2023tree}, or trajectories \citep{snell2024scaling} to accumulate broader candidate evidence for subsequent reasoning. 
Despite their empirical success in adapting to particular reasoning strategies, both lines impose reasoning control from outside the model (i.e., through pre-specified formats or search policies) and ignore its evolving internal state as a signal for what evidence is actually needed at subsequent decisions, limiting generalization across heterogeneous reasoning structures \citep{hassid2025dont} and efficiency under constrained inference \citep{zhou2023least}.

% In contrast, we argue that the model's \emph{endogenous state} (i.e., the internally evolving state induced by its own reasoning trajectory) is what genuinely organizes the evidence context needed for subsequent decisions.
Prior work has shown that LLMs internals exhibit structured representation geometry \citep{zou2023repe}, trajectory-level regularities \citep{sun2026llm}, and confidence dynamics \citep{tikhonov2026confidence}. 
However, these signals are typically used descriptively, such as monitoring and reasoning-trajectory interpretation \citep{ballon2026probing,damirchi2026truth,feng2025monitoring}, or early stopping criteria \citep{hosseini2026early}, rather than to govern reasoning itself.
Together, these observations point to the significance of the model's \emph{endogenous state} (i.e., the internally evolving state induced by its own reasoning trajectory).
In parallel, reasoning also depends critically on whether decisions are steadily grounded in suitable and decision-relevant historical evidence rather than generic or noisy history \citep{creswell2022selection,press2023measuring,wang2024rat}. Existing approaches address this through context retrieval and compression \citep{jiang2024longllmlingua,li2023compressing,packer2023memgpt}, or latent-space reasoning \citep{hao2025training,wang2025system}, both of which bypass endogenous state as a vital signal for organizing the evidence genuinely needed for subsequent steps.
Appendix~\ref{app:related_work} provides further related discussion.
% Together, these observations point to a reasoning paradigm that uses structured endogenous states to dynamically control and organize evidence context with high quality for subsequent reasoning.

This motivates a reframing of reasoning itself: \emph{effective reasoning is actually a dynamic process of organizing the right evidence context for the next decision under the current endogenous state, rather than following an externally prescribed token chain.} 
We introduce \textbf{State of Thought (SoT)}, a reasoning paradigm for \emph{endogenous reasoning}. 
At each reasoning step, SoT reads out a compact dynamics-geometric state from the model's internal information transfer and uses it to selectively activate the historical evidence context most useful for the next decision.
The active reasoning context is therefore not generated by pre-specified reasoning programs or broader searches, but by the model's own evolving state. 
In this way, SoT turns reasoning into a closed loop: endogenous state organizes active evidence, which conditions the next reasoning step, and the new step in turn updates the state.

We evaluate SoT at scale on three LLM backbones over 16 datasets and one multimodal backbone over 4 datasets. 
Across the 16-dataset LLM suite, SoT improves accuracy by \(85.0\%\), \(59.1\%\), and \(49.6\%\) over baselines on Llama-3.1-8B, Qwen2.5-14B, and Mixtral-8x7B, while reducing tokens/latency by \(40.8\%\)/\(54.6\%\), \(84.5\%\)/\(68.6\%\), and \(88.5\%\)/\(14.7\%\), respectively.
On Qwen3-VL-8B, it further improves accuracy by \(7.8\%\) across 4 multimodal benchmarks.
As exploratory extensions under constrained access, a training-free controller and an embedding-only controller still achieve \(38.6\%\) and \(36.1\%\) improvement over baselines on accuracy across models, while a trajectory-level post-hoc judge reaches \(84.1\%\) agreement on three API-accessed models.
Together, these results suggest that changing the control variable of reasoning—from external token programs to endogenous state-conditioned evidence organization—offers a scalable and robust route to reasoning improvement.


% We summarize our contributions as follows.
% \begin{itemize}
%     \item We first formulate test-time reasoning as endogenous state-conditioned evidence organization. This view reframes reasoning control from externally prescribed procedures to the model's evolving reasoning state.

%    \item We introduce State of Thought, a lightweight controller that reads a compact dynamics-geometric state from the model and uses it to coordinate sparse evidence activation with reasoning continuation.

%     \item we evaluate SoT across arithmetic, knowledge-intensive, symbolic, code, long-context, and multimodal reasoning settings. The results show consistent gains in accuracy and efficiency over scripted, sampling-based, search-based, memory-oriented, latent-reasoning, and RL-based baselines.
% \end{itemize}